\documentclass{article}
\usepackage{preamble}

\begin{document}
\header{Filemon Mateus}{CS 6140}{Dimensionality Reduction}{Homework \# 7}{\today}
\begin{enumerate}[leftmargin={*}, font={\bf}, label={\arabic*.}, ref={\arabic*}]
  \item \label{qst:1}
    \begin{enumerate}[label={(\roman*)}, ref={\theenumi(\roman*)}]
      \item \label{qst:1i}
        The following is a scatter plot consisting of $186$ data points along with
        two principal directions.
        \begin{figure}[H]
          \centering
          \begin{tikzpicture}
            \begin{axis}[
              xlabel={Child Mortality},
              ylabel={Women Fertility},
              clip marker paths={true}
            ]
              \addplot[blue, only marks, mark size=1pt] table[col sep=comma] {dev-dir/week7/data/mortality-vs-fertility.csv};
              \addplot[-stealth, line width=2pt, red] coordinates {(0, 0) (-0.93713505 * 3, -0.37473738 * 3)};
              \addplot[-stealth, line width=2pt, green!60!black] coordinates {(0, 0) (0.37473738 * 1, -0.93713505 * 1)};
            \end{axis}
          \end{tikzpicture}
          \caption{
            Women fertility as a function of neonatal child mortality.
          }
          \label{fig:mortality-vs-fertility}
        \end{figure}
        Based on \autoref{fig:mortality-vs-fertility}, the direction captured by
        the vector in green (along the coordinates $\left[0.3747, -0.9271\right]$)
        carries the maximal amount of variance across the $186$ observations, and
        thus, it encodes more information to that of the vector/direction in red.
        It is therefore reasonable to conclude that the direction of the vector in
        green is more important.
        
      \item \label{qst:1ii}
        The following is an approximation of the mortality/fertility dataset with
        the best Rank-$1$ approximation.
        \begin{figure}[H]
          \centering
          \begin{tikzpicture}
            \begin{axis}[
              xlabel={Child Mortality},
              ylabel={Women Fertility},
              clip marker paths={true}
            ]
              \addplot[blue, only marks, mark size=1pt] table[col sep=comma] {dev-dir/week7/data/mortality-vs-fertility.csv};
              \addlegendentry{Ground Truth Observations};
              \addplot[red, only marks, mark size=1pt] table[col sep=comma] {dev-dir/week7/data/mortality-vs-fertility-rank-1-approx.csv};
              \addlegendentry{Best Rank-$1$ Approximation};
            \end{axis}
          \end{tikzpicture}
          \caption{
            Women fertility as a function of neonatal child mortality. \\
            Best Rank-$1$ approximation.
          }
        \end{figure}
    \end{enumerate}

  \item \label{qst:2}
    \begin{enumerate}[label={(\Alph*)}, ref={\theenumi(\Alph*)}]
      \item \label{qst:2A}
        \begin{enumerate}[label={(\roman*)}, ref={\theenumii(\roman*)}]
          \item
            Below is a plot with $\varepsilon$ on the $x$-axis and the mean absolute
            differences on the $y$-axis corresponding to the gaussian random projection
            subroutine.
            \begin{figure}[H]
              \centering
              \begin{tikzpicture}
                \begin{axis}[
                  xlabel={$\varepsilon$-values},
                  ylabel={Mean Absolute Difference},
                  scaled ticks={false},
                  yticklabel style={
                    /pgf/number format/fixed,
                    /pgf/number format/zerofill,
                    /pgf/number format/precision=2
                  },
                ]
                  \addplot[
                    blue,
                    thick,
                    mark=*
                  ] table[col sep=comma] {dev-dir/week7/data/gauss-proj-eps-vs-mean-abs-diff.csv};
                \end{axis}
              \end{tikzpicture}
              \caption{
                Mean absolute difference of gaussian random projection.
              }
            \end{figure}
            
          \item
            Below is a histogram corresponding to the absolute differences.
            \begin{figure}[H]
              \centering
              \begin{tikzpicture}
                \begin{axis}[
                  xlabel={Absolute Difference},
                  ylabel={Frequency},
                  scaled ticks={false},
                  xticklabel style={
                    /pgf/number format/fixed,
                    /pgf/number format/zerofill,
                    /pgf/number format/precision=2
                  },
                  yticklabel style={
                    /pgf/number format/fixed
                  },
                  ymin={0}
                ]
                  \addplot [
                    fill={blue},
                    draw={black},
                    hist={data=x, bins=10}
                  ] file {dev-dir/week7/data/gauss-proj-abs-diff.csv};
                \end{axis}
              \end{tikzpicture}
              \caption{
                Histogram of absolute differences of gaussian random projection.
              }
            \end{figure}
            
          \item
            Yes, from mean absolute difference plot and absolute differences histogram,
            we see that pairwise distances are nearly preserved.

          \item
            Below is the implementation:
            \lstinputlisting
              [caption={\sf gauss-rand-proj.py}, label={lst:gauss-rand-proj}]
              {dev-dir/week7/gauss-rand-proj.py}
        \end{enumerate}
        
      \item \label{qst:2B}
        \begin{enumerate}[label={(\roman*)}, ref={\theenumii(\roman*)}]
          \item
            Below is a plot with $\varepsilon$ on the $x$-axis and the mean absolute
            differences on the $y$-axis corresponding to the sparse random projection
            subroutine.
            \begin{figure}[H]
              \centering
              \begin{tikzpicture}
                \begin{axis}[
                  xlabel={$\varepsilon$-values},
                  ylabel={Mean Absolute Difference},
                  scaled ticks={false},
                  yticklabel style={
                    /pgf/number format/fixed,
                    /pgf/number format/zerofill,
                    /pgf/number format/precision=2
                  },
                ]
                  \addplot[
                    thick,
                    blue,
                    mark=*
                  ] table[col sep=comma] {dev-dir/week7/data/sparse-proj-eps-vs-mean-abs-diff.csv};
                \end{axis}
              \end{tikzpicture}
              \caption{
                Mean absolute difference of sparse random projection.
              }
            \end{figure}
             
          \item
            Below is a histogram corresponding to the absolute differences.
            \begin{figure}[H]
              \centering
              \begin{tikzpicture}
                \begin{axis}[
                  xlabel={Absolute Difference},
                  ylabel={Frequency},
                  scaled ticks={false},
                  xticklabel style={
                    /pgf/number format/fixed,
                    /pgf/number format/zerofill,
                    /pgf/number format/precision=2
                  },
                  yticklabel style={
                    /pgf/number format/fixed
                  },
                  ymin={0}
                ]
                  \addplot [
                    fill={blue},
                    draw={black},
                    hist={data=x, bins=10}
                  ] file {dev-dir/week7/data/sparse-proj-abs-diff.csv};
                \end{axis}
              \end{tikzpicture}
              \caption{
                Histogram of absolute differences of sparse random projection.
              }
            \end{figure}
            
          \item
            Yes, from the mean absolute difference plot and absolute differences histogram,
            we see that pairwise distances are nearly preserved.
            
          \item
            Below is the implementation:
            \lstinputlisting
              [caption={\sf sparse-rand-proj.py}, label={lst:sparse-rand-proj}]
              {dev-dir/week7/sparse-rand-proj.py}
        \end{enumerate}
    \end{enumerate}
\end{enumerate}
\end{document}
